\section{Cyclic shifts and the relative central kernel}\label{sec:shifts}

We now study elements of $J(R)$ whose reduction to $\St(\F_2)$ is
trivial.  The grading lets us evaluate such an element at a finite
numerical matrix.  We shall prove that evaluation at a sufficiently
long cyclic shift is trivial, and that evaluation at a cyclic shift
of $2$-power size is a power of the original element.  These two
identities imply $2$-primary torsion in the relative central kernel.
Later we will compare the first identity with an odd power.

The matrix substitutions are elementary versions of the graded
transfer operations of Weibel \cite[Section~1]{Weibel1987}.
At the level of stable algebraic $K$-theory, the relative torsion
conclusion also follows from his continuous Witt-vector action
\cite[Theorem~0.1]{Weibel1987}.
We reproduce the concrete Steinberg-group arguments of the companion
\cite[Lemmas~6.4--6.7]{Companion}, since their exact identities, rather
than identities of matrix images, are needed here.

\subsection{Block roots and triangular evaluation}

For the next three lemmas, let $B$ be any unital $\F_2$-algebra.
Fix $n\geq3$ and $b\geq1$, and label $nb$ coordinates by
$(i,s)$, where $1\leq i\leq n$ and $0\leq s<b$.
We call $i$ the outer index and $s$ the inner index.  For
$A=(a_{st})\in\Mat_b(B)$ and $i\ne j$, define the \emph{block root}
\[
 X_{ij}(A)=\prod_{s,t=0}^{b-1}
           x_{(i,s),(j,t)}(a_{st})\in\St_{nb}(B).
\]
All factors commute, because their source outer index is $i$ and
their target outer index is the different index $j$.

\begin{lemma}[Block roots and coordinate changes]\label{lem:block-laws}
The block roots satisfy the Steinberg relations over $\Mat_b(B)$:
\begin{align*}
 X_{ij}(A+A')&=X_{ij}(A)X_{ij}(A'),\\
 [X_{ij}(A),X_{kl}(A')]&=1
       &&\text{if }j\ne k\text{ and }i\ne l,\\
 [X_{ij}(A),X_{jk}(A')]&=X_{ik}(AA')
       &&\text{if }i,j,k\text{ are distinct}.
\end{align*}
Consequently they define a homomorphism
$\St_n(\Mat_b(B))\longrightarrow\St_{nb}(B)$.
Moreover, for every $P\in\GL_b(\F_2)$ there is an element
$q\in\St_{nb}(B)$ such that
\[
 q^{-1}X_{ij}(A)q=X_{ij}(P^{-1}AP)
 \qquad(i\ne j).
\]
Its matrix image is the block diagonal matrix with a copy of $P$
in every outer block.
\end{lemma}

\begin{proof}
Additivity and the commuting relation follow factor by factor.
For the consecutive commutator, the only nontrivial elementary
commutators are
\[
 [x_{(i,s),(j,t)}(a_{st}),
   x_{(j,t),(k,v)}(a'_{tv})]
   =x_{(i,s),(k,v)}(a_{st}a'_{tv}).
\]
Each resulting root commutes with both input blocks.  Expanding the
commutator and using additivity at each position therefore gives
the entries $\sum_t a_{st}a'_{tv}$ of $AA'$.

An elementary operation supported inside outer block $i$ acts on
the outgoing strip to block $j$ by the corresponding row operation;
an operation inside block $j$ acts by the inverse column operation.
These are exact strip actions by Lemma~\ref{lem:strip}.  An operation
inside any other outer block centralizes the block root.  Gaussian
elimination over $\F_2$ writes $P$ as a product of elementary
matrices: every nonzero pivot is $1$, and a row transposition is
a product of three row additions.  Lift those factors inside each
outer block.  Applying the strip actions successively gives the
asserted conjugation.  For $b=1$, $P=1$ and the assertion is immediate.
\end{proof}

If $v(t)$ is a root word over the polynomial algebra $B[t]$, with
$t$ central, and $T\in\Mat_b(\F_2)$, let $W_v(T)$ be its evaluation
at $t=T$, followed by the block-root homomorphism.  The substitution
$B[t]\to\Mat_b(B)$ is a unital algebra homomorphism: the numerical
matrix $T$ commutes with every scalar matrix from $B$.
In particular, if $v(t)$ has identity matrix image over $B[t]$,
then $W_v(T)$ has identity matrix image over $B$.
For each inner index $s$, also write
\[
 \iota_s:\St_n(B)\longrightarrow\St_{nb}(B),\qquad
 x_{ij}(a)\longmapsto x_{(i,s),(j,s)}(a).
\]
These coordinate homomorphisms have commuting images for distinct
$s$; we do not assume they are injective.

\begin{lemma}[Triangular evaluation]\label{lem:triangular}
Let $U$ be generated by the roots
$x_{(i,s),(j,t)}(a)$ with $s<t$, allowing $i=j$.
Then $U$ injects into its matrix image and is normalized by every
root whose inner indices agree.  If a root word $v(t)$ over $B[t]$
has identity matrix image, and $T\in\Mat_b(\F_2)$ is upper
triangular with constant diagonal entry $\lambda$, then
\begin{equation}\label{eq:triangular-evaluation}
 W_v(T)=\prod_{s=0}^{b-1}\iota_s\bigl(v(\lambda)\bigr).
\end{equation}
\end{lemma}

\begin{proof}
The \emph{gap} of a generator of $U$ is $t-s>0$.
Two such roots either commute or have a commutator root whose
gap is the sum of their gaps, by the consecutive-root relations
in the two possible orders.  They cannot be opposite roots.
For $g\geq1$, let $U_g$ be generated by roots of gap at least $g$;
thus $U_1=U$ and $U_b=1$.  Conjugating a generator of $U_g$ by
a generator of $U$ adds only roots of larger gap, so $U_g$ is
normal in $U$.  Modulo $U_{g+1}$, the group $U_g$ is generated
by the commuting roots of gap $g$.

It follows, successively for $g=1,\ldots,b-1$, that every element
of $U$ is an ordered product with one factor at each coordinate
pair, in increasing gap; repeated coefficients at a given pair
are added.  This is a finite collection process, since there are
finitely many coordinate pairs and $U_b=1$.
If its matrix image is the identity, the entries of gap one
are precisely the corresponding collected coefficients, so they
vanish.  After coefficients of smaller gap have vanished, the
same assertion holds at gap $g$: any product of two remaining
entry matrices has gap greater than $g$.  Induction makes every
coefficient zero and proves faithfulness.

A root with equal inner indices cannot be opposite to a
positive-gap root.  Its commutator with a positive-gap root is
either trivial or has the same positive gap.  This proves that
it normalizes $U$.

For a polynomial entry $a(t)$, the matrix $a(T)$ is upper
triangular with diagonal entries $a(\lambda)$.  Split its block
root into its commuting diagonal and strictly upper factors.
The diagonal factors normalize $U$, so the strictly upper
factors in the complete word can be collected on the right.
Diagonal factors at different inner indices commute, giving
\[
 W_v(T)=\left(\prod_s\iota_s(v(\lambda))\right)u,
 \qquad u\in U.
\]
The left side and all the words $v(\lambda)$ have identity
matrix image.  Hence $u$ has identity matrix image.  Faithfulness
of $U$ gives $u=1$ and proves \eqref{eq:triangular-evaluation}.
\end{proof}

\subsection{Removing the wrap of a cyclic shift}

Let $C_b$ be the numerical cyclic shift from row $s$ to column
$s+1$ modulo $b$, and let $S_b$ be the shift with the same entries
except for the entry from row $b-1$ to column $0$, which is removed.
Thus $S_b$ is strictly upper triangular.  The next lemma replaces
$C_b$ by $S_b$ without changing the evaluated Steinberg element.
Its proof confines the difference to a faithful strip between
two disjoint bands of inner indices.

\begin{lemma}[Removing the wrap]\label{lem:wrap}
Let $v(t)$ be a product of $m\geq1$ root factors over $B[t]$,
whose polynomial entries have degrees at most $d\geq0$.
Suppose its matrix image is the identity over $B[t]$.
Put $\delta=\max\{1,d\}$.  If
\[
 b>2(m+1)\delta,
\]
then $W_v(C_b)=W_v(S_b)$ in $\St_{nb}(B)$.
\end{lemma}

\begin{proof}
For each factor, split its cyclic evaluation as $D_hF_h$, where
$D_h$ is its evaluation at $S_b$ and $F_h$ consists of the entries
that wrap past inner index $b-1$.  The splitting is exact because
all roots in one outer rectangle commute.  Since $d<b$, a term
of degree $a\leq d$ goes from $s$ to $s+a$ and can wrap only once.
Every root in $F_h$ therefore has source in
$\{b-\delta,\ldots,b-1\}$ and target in
$\{0,\ldots,\delta-1\}$.

For $r\geq1$, set
\[
 H_r=\{b-r\delta,\ldots,b-1\},\qquad
 L_r=\{0,\ldots,r\delta-1\}.
\]
When $b>2r\delta$, the roots from inner indices in $H_r$ to
inner indices in $L_r$, with arbitrary outer indices, form an
abelian strip $V_r$ that injects into matrices by
Lemma~\ref{lem:strip}.  Figure~\ref{fig:wrap-strip} displays these
two bands.  We shall show that conjugation by one $D_h$ enlarges
their widths by at most $\delta$.

\begin{figure}[htbp]
\centering
\begin{tikzpicture}[x=0.95cm,y=0.8cm,>=Stealth]
 \fill[gray!15] (0,0) rectangle (2,0.65);
 \fill[gray!30] (8,0) rectangle (10,0.65);
 \draw (0,0) rectangle (10,0.65);
 \draw (2,0)--(2,0.65) (8,0)--(8,0.65);
 \node at (1,0.325) {$L_r$};
 \node at (9,0.325) {$H_r$};
 \draw[->] (9,0.85)--(9,1.4)--(1,1.4)--(1,0.85);
 \node at (5,1.77) {correction roots};
 \node[below] at (0,0) {$0$};
 \node[below] at (2,0) {$r\delta-1$};
 \node[below] at (8,0) {$b-r\delta$};
 \node[below] at (10,0) {$b-1$};
\end{tikzpicture}
\caption{The correction strip $V_r$ goes from the right band
$H_r$ to the left band $L_r$; each band includes every outer index.
The proof shows that one truncated block factor moves each endpoint
by at most $\delta$.  The bound $b>2(m+1)\delta$ keeps the enlarged
bands disjoint throughout the collection.}
\label{fig:wrap-strip}
\end{figure}

More precisely, for $r\leq m$ we claim that
\begin{equation}\label{eq:strip-enlargement}
 D_h^{-1}V_rD_h\subseteq V_{r+1}.
\end{equation}
Fix a root of $V_r$.  The constituents of $D_h$ all have the
same outer source $i$ and outer target $j\ne i$, and inner
displacements between $0$ and $d$.  They commute with one
another.  A constituent that commutes with the initial root
also commutes with every conjugate of it by the other
constituents, and so can be omitted from this conjugation.

An interacting constituent changes the source endpoint only
when its target is the old source; the new source is at most
$d$ positions to the left.  It changes the target endpoint only
when its source is the old target; the new target is at most
$d$ positions to the right.  These statements follow directly
from the consecutive-root relations, including their reverse
order.  An opposite-root interaction cannot occur.  Initially
the source inner index minus the target inner index is greater
than $b-2r\delta$.  After one endpoint has moved, the difference
is still greater than
\[
 b-(2r+1)\delta>\delta\geq d.
\]
Thus no constituent, whose forward displacement is at most $d$,
can be opposite to either the initial root or a root obtained
after one endpoint move.

Nor can displacements accumulate repeatedly at one endpoint
during this conjugation.  A changed source has outer index $i$,
whereas every constituent has target outer index $j\ne i$;
the source cannot change again.  A changed target has outer
index $j$, whereas every constituent has source outer index
$i$; the target cannot change again.  Both endpoints may change,
but each changes at most once.  A root with both endpoints
changed has outer indices $i,j$ and commutes with every
constituent.  All roots produced lie in $V_{r+1}$, proving
\eqref{eq:strip-enlargement}.  This reasoning applies to a product
of strip roots as well.

It remains to collect the corrections without conjugating
them by one another.  Put $D_{>h}=D_{h+1}\cdots D_m$ and
$D_{>m}=1$.  The exact group identity
\[
 \prod_{h=1}^m(D_hF_h)
  =(D_1\cdots D_m)
       \prod_{h=1}^m(D_{>h}^{-1}F_hD_{>h}),
\]
where the final product is in increasing order of $h$, does this.
Every $F_h$ starts in $V_1$, and every suffix has at most $m$
factors.  By \eqref{eq:strip-enlargement}, all final corrections
belong to $V_{m+1}$.  Therefore
\[
 W_v(C_b)=W_v(S_b)u,\qquad u\in V_{m+1}.
\]
Both evaluated words have identity matrix image.  The same is
true of $u$, and the disjointness of $H_{m+1}$ and $L_{m+1}$
makes $V_{m+1}$ faithful.  Hence $u=1$.
\end{proof}

\subsection{Graded evaluation and powers of central elements}

Return to the graded algebra $R=\bigoplus_{a\geq0}R_a$ with
$R_0=\F_2$.  The map
\[
 \gamma:R\longrightarrow R[t],\qquad
       \sum_a r_a\longmapsto\sum_a r_a t^a
       \quad(r_a\in R_a)
\]
is a unital algebra homomorphism.  Applying it to a root word
representing $z\in J(R)$ gives a word $z(t)$ with identity
matrix image over $R[t]$.  Its evaluations at $1$ and $0$ are
$z$ and the degree-zero reduction of $z$, respectively.

For $T\in\Mat_b(\F_2)$, fix an injection from the pairs $(i,s)$,
where $i$ ranges over the stable index set and $0\leq s<b$,
into the stable index set.  Define $W_z(T)$ by applying the
graded substitution and evaluating by block roots in those
coordinates.  This is a homomorphism on $\St(R)$ by
Lemma~\ref{lem:block-laws}, so $W_z(T)$ depends only on $z$.
The element $W_z(T)$ belongs to $J(R)$.  All words and their
block evaluations use finitely many coordinates; choosing a
different injection gives the same element by centrality and
Lemma~\ref{lem:permutations}.

\begin{lemma}[Vanishing at long cyclic shifts]\label{lem:cyclic-vanishing}
If $z\in J(R)$ maps to $1$ in $\St(\F_2)$, then
\[
 W_z(C_b)=1
 \qquad\text{for every sufficiently large integer }b.
\]
More precisely, if the graded word $z(t)$ has $m\geq1$ factors
of degrees at most $d$, it suffices that
$b>2(m+1)\max\{1,d\}$.
\end{lemma}

\begin{proof}
Choose one finite root word for $z$, and apply
Lemma~\ref{lem:wrap} to its graded word.  Since $S_b$ is strictly
upper triangular, Lemma~\ref{lem:triangular} gives, in the stable
group,
\[
 W_z(C_b)=W_z(S_b)
       =\prod_{s=0}^{b-1}\iota_s(z(0))=1.
\]
The fixed injection extends each $\iota_s$ to a homomorphism on
the stable group.  Since $z(0)$ is trivial there by hypothesis,
every displayed coordinate image is trivial.  The empty word
satisfies the conclusion for all $b$.
\end{proof}

\begin{proposition}[Relative $2$-primary torsion]\label{prop:relative-torsion}
Every element of
\[
 J(R)\cap\ker\bigl(\St(R)\longrightarrow\St(\F_2)\bigr)
\]
has finite $2$-power order.
\end{proposition}

\begin{proof}
Let $z$ be such an element.  If $b$ is a power of two, then
$C_b^b=\Id$ and characteristic two gives
$(C_b-\Id)^b=0$.  A basis adapted to the successive kernels of
the nilpotent map $C_b-\Id$ makes $C_b$ upper triangular with
all diagonal entries $1$.  Thus some
$P\in\GL_b(\F_2)$ has this effect.  By
Lemma~\ref{lem:block-laws}, conjugation by a Steinberg lift of
the repeated basis change takes $W_z(C_b)$ to
$W_z(P^{-1}C_bP)$.  Centrality of $J(R)$ makes these elements
equal.  Lemma~\ref{lem:triangular} then gives
\begin{equation}\label{eq:cyclic-two-power}
 W_z(C_b)=\prod_{s=0}^{b-1}\iota_s(z)=z^b.
\end{equation}
The last equality uses Lemma~\ref{lem:permutations}: every
coordinate image of the central element $z$ equals $z$ in the
stable group.  Choose a power of two $b$ large enough for
Lemma~\ref{lem:cyclic-vanishing}.  Then $z^b=1$.
\end{proof}

Only the positive grading and the degree-zero field were used
in this section.  To detect central elements by finite degree
truncations, we next use the additional phase structure of the
specific algebra $R$.  It will yield an odd-power counterpart
of \eqref{eq:cyclic-two-power}.
